\chapter{KESIMPULAN DAN SARAN}

% TODO: Isi bab ini setelah eksperimen dan analisis selesai. Kerangka di
% bawah disusun berdasarkan dimensi evaluasi pada Bab I dan analisis Bab IV.

\section{Kesimpulan}

% TODO: Rumuskan kesimpulan untuk setiap dimensi evaluasi:
% - Konsistensi deployment: bagaimana perbedaan antara paradigma manual dan
%   GitOps setelah skenario penyimpangan (skenario A--G)?
% - Waktu deteksi, pemulihan, dan identifikasi akar masalah: seberapa besar
%   perbedaan antara lingkungan otomatis (ArgoCD) dan intervensi manual?
% - Traceability operasional: sejauh mana perbedaan kelengkapan jejak
%   perubahan antara Git (Lingkungan B) dan audit logs (Lingkungan A)?
% Sertakan pula kesimpulan utama mengenai skenario penyimpangan yang paling
% berdampak dan rekomendasi pemilihan paradigma deployment.

\section{Saran}

% TODO: Rumuskan saran untuk penelitian lanjutan:
% - Replikasi dengan operator berbeda tingkat keahlian (mitigasi batasan
%   operator tunggal).
% - Evaluasi paradigma perantara dan push-based pada spektrum
%   deklaratif-imperatif.
% - Perluasan ke klaster multi-cluster, workload lain, atau tool GitOps lain
%   (mis. FluxCD).
% - Saran praktis bagi praktisi industri terkait kondisi di mana manual
%   deployment tetap memadai.

\section{Keterbatasan Penelitian}

% TODO: Dokumentasikan keterbatasan penelitian yang teridentifikasi selama
% pelaksanaan eksperimen, mengacu pada ancaman validitas pada Bab III
% (validitas internal, eksternal, konstruk, dan reliabilitas), serta
% kendala operasional yang muncul selama pengumpulan data.
